% HD-R15. GPT-6 Astra Pro /root, 21 November 2026.
\subsection{HD-R15: terminating reduction through actual coefficient blocks}
\label{sec:rank-three-source-reduction}

Consider an actual natural finite-curve source with period
rank nine for a sixfold component still missing in HD-S43:
over \(K=\mathbb Q(i)\) its ordinary determinant is outside
\([-1],[-3]\), and over
\(K=\mathbb Q(\sqrt{-d})\), \(d=2,7,11,19,43,67,163\),
it is outside \([-1]\). The coefficient representation is
faithful and absolutely irreducible, with exact character
field \(K\), and the original primitive Jacobian corner
and its polarization are prescribed.

There is a finite chain of actual intermediate-curve
quotients and integer-cleared primitive-corner isogenies
that carries this source to one of the four forms below.
The final coefficient rank is no larger than the original
rank. The original polarized sixfold is retained up to the
displayed actual constant isogenies; its period rank and
ordinary determinant class are unchanged.
\begin{enumerate}
\item Every nonidentity coefficient element has moving
rank at least four. Its coefficient rank is at least~21.
\item The subgroup generated by the actual rank-three
elements is absolutely irreducible and complex primitive.
Its coefficient rank lies in \(15\le m\le35\), with the
field-dependent intervals and rank-three spectra of HD-G31.
\item That rank-three subgroup is absolutely irreducible
and complex imprimitive. Every nontrivial block system
has block dimension at most three. Its coefficient rank
is at least~15.
\item The coefficient group and representation are
literally \(S_s\times\Gamma\) on
\(\operatorname{Aug}(K^s)\otimes_K M\), with
\(s=r+1\), \(r\ge9\), and \(\dim_KM=3\), as in
HD-S45. The finite absolutely irreducible group
\(\Gamma\subset GL_K(M)\) has exact character field
\(K\) and has no \(K\)-defined line system.
\end{enumerate}
These are reductions of actual sources, without an assertion
that any surviving form gives a source or a new detector.
The field-specific lower ranks of HD-R12 remain in force
in each form. In form~2 they exclude \(d=163\) entirely.

\paragraph{1. One descended block gives a strictly smaller actual source.}
Let \(\rho:G\to GL_K(V)\) be the current representation,
and suppose its actual orthogonal \(K\)-block system is
\(V=\bigoplus_{i\in I}T_i\), with \(|I|>1\), permuted
transitively by \(G\). Choose \(T=T_i\), let \(J\) be
its stabilizer, and let \(\tau:J\to GL_K(T)\) be the
restricted action, with kernel \(J_0\). The actual induction
isomorphism and its inverse are
\[
 \begin{aligned}
 K[G]\otimes_{K[J]}T&\longrightarrow V,
                 &&g\otimes t\longmapsto\rho(g)t,\\
 V&\longrightarrow K[G]\otimes_{K[J]}T,
                 &&v\longmapsto
      \sum_{gJ\in G/J}g\otimes\pi_T\rho(g^{-1})v.
 \end{aligned}                                              \tag{R15.1}
\]
Orthogonality of the block projectors verifies both
composites. The coset relation verifies independence of
the chosen representatives. A proper \(J\)-invariant
complex subspace of \(T_{\mathbb C}\) would give a proper
\(G\)-invariant direct sum through (R15.1). Thus the
faithful quotient \(J/J_0\) is absolutely irreducible
on \(T\).

Its character is \(K\)-valued. The equality
\(\chi_\rho=\operatorname{Ind}_J^G\chi_\tau\) makes every
value of \(\chi_\rho\) an integral sum of values of
\(\chi_\tau\). Since the first exact field is \(K\),
the field of the second character both contains and is
contained in \(K\); it is exactly \(K\). This argument
also applies after replacing \(J\) by its faithful image.

For the original connected cover \(\widetilde C\), use
the actual quotient curves and map
\[
       Y=\widetilde C/J,\qquad
       C'=\widetilde C/J_0,\qquad
       \alpha:\widetilde C\longrightarrow C'.             \tag{R15.2}
\]
To fix the local-system convention, write the induced
functions as \(F_v(x)=\pi_T\rho(x)v\), satisfying
\(F(hx)=\tau(h)F(x)\); their inverse is
\(\sum_{J\backslash G}\rho(x)^{-1}F(x)\).
These same coefficient maps commute with taking inertia
invariants, so extend across every branch stalk. They
identify the parabolic cohomology and its Hodge filtration
through (R15.2), including every branch value that is
unramified for the quotient coefficient representation.

Choose \(0\ne u\in T(K)\). Its normalized rational
matrix-coefficient idempotent is the same element when
computed in \(\mathbb Q[G]\) and in \(\mathbb Q[J]\):
off \(J\) its coefficients vanish by orthogonality, and
\(\dim V/|G|=\dim T/|J|\). The quotient norm equations
give the actual polarized pullback on the primitive factors,
\[
        \alpha_*\alpha^*=[|J_0|],\qquad
        (\alpha^*)^*\theta_{J\widetilde C}
                            =|J_0|\theta_{JC'}.            \tag{R15.3}
\]
Let \(w\) be the originally prescribed primitive vector.
To retain that corner, use the exact Fourier operator
\(b_{w,u}\) of (S41.15)--(S41.16), and choose an integer
\(M\) clearing its coefficients. It has
\[
       b_{w,u}^{\dagger}b_{w,u}
          =\frac{h(w,w)}{h(u,u)}e_u,
 \qquad
 (Mb_{w,u})^*\theta_{A_w}
          =M^2\frac{h(w,w)}{h(u,u)}\theta_{A_u}.            \tag{R15.4}
\]
Its composition with the actual quotient pullback is an
isogeny from the new primitive factor to the prescribed
one. A rational inverse is given by the norm and the
adjoint divided by their displayed nonzero scalars.

Every multiplier is a positive rational number. For the
six-dimensional Hermitian multiplicity space its sixth
power is \(N_{K/\mathbb Q}(t^3)\). Thus the ordinary
determinant is unchanged. Over a finite parameter cover
fixing the group labels and coefficient lattices, these
are actual fixed isogenies of families; their maps of
Hodge filtrations identify the period differentials.
The new coefficient rank is \(\dim T=\dim V/|I|<\dim V\).

\paragraph{2. Applying the rank-three support decomposition.}
Every current source still has a missing determinant.
Rank-one sources are covered by the actual cyclic
quotients in HD-S34; over the seven fields with only
\(\{\pm1\}\) as roots of unity their character field
would instead be rational. Rank-two sources have period
rank at most seven by HD-S30. Thus the current coefficient
rank is at least three. In genus one the Euler formula
allows at most six true coefficient branch values and
hence at most six parameters. In genus at least two it
forces coefficient rank three, genus two and no branch,
whose base moduli dimension is three. Consequently period
rank nine forces a sphere base. HD-S40 now excludes
coefficient ranks at most14. HD-S43 and its cited
reflection and bireflection proofs also exclude every
actual coefficient element of moving rank one or two.
Thus the current rank is at least15, and its sphere tuple
has at least twelve branch values, each of codimension
at least three.

Form its actual normal subgroup \(N_3\) generated by all
moving-rank-three elements. For \(N_3\ne1\), HD-S44
constructs the complete isotypic projectors of \(K[N_3]\)
and gives the following alternatives. At multiplicity one,
the projectors descend individually, or the Galois action
has cyclic cubic orbits. The latter alternative is excluded
by S44's actual semilinear quotient maps or its low-rank
complete-block reduction. At multiplicity three the
projectors descend and have the explicit
\(AP\otimes_K PU\) coordinates.

In either descended case, more than one isotypic component
gives the proper \(K\)-block system of paragraph1.
Its projectors are orthogonal for the prescribed form,
by the character projector's adjoint formula in HD-S44.
Pass to the strictly smaller source (R15.2)--(R15.4) and
repeat the same analysis on its actual faithful group.
This procedure terminates, because at each repetition its
positive integer coefficient rank strictly decreases.
Only the finitely many actual groups, quotients and
integer-cleared maps encountered in that procedure are
used; no limit of representations is taken.

At termination, \(N_3\) is trivial, or it has one
isotypic component. Multiplicity one in the latter case
means that \(N_3\) is absolutely irreducible.
Multiplicity three has exactly the tensor system treated
by HD-S45, whose remaining actual product is form4.
Its \(K\)-defined line-system case is covered by that
proof, so it is removed from the retained form.

\paragraph{3. Exact ranks and the imprimitive block bound.}
For trivial \(N_3\), every true inertia has codimension
at least four. The original parabolic dimension and
pointed-sphere parameter count give
\[
             2m+6=\sum_{j=1}^{b}c_j\ge4b\ge48,
                          \qquad m\ge21.                  \tag{R15.5}
\]
This is form1. For absolutely irreducible \(N_3\), HD-R14
excludes every tensor structure on its actual
complexification. Its complex-primitive branch is exactly
HD-G31, yielding form2 and its precise field intervals.

For completeness, let \(\bigoplus_{i=1}^a V_i\) be
any nontrivial system of imprimitivity for the remaining
absolutely irreducible \(N_3\), of common block dimension
\(b\). Some rank-three generator must permute its blocks
nontrivially; otherwise each block would be invariant under
all generators. On a moved cycle of length \(\ell\ge2\),
write its actual maps successively as
\(A_i:V_i\to V_{i+1}\), with the last returning to the
first. A fixed vector is determined by its first entry,
which must lie in
\(\ker(A_\ell\cdots A_1-I)\). The literal map sending
that entry to its successive images identifies the fixed
space with this kernel. Its dimension is at most \(b\),
so this cycle moves at least \(b(\ell-1)\) dimensions.
The total moving rank is three, and therefore \(b\le3\).
This proves form3 without asserting that this complex
block system descends or is preserved by an overgroup.

The chain (R15.1)--(R15.4) proves every claimed comparison
of the final form with the specified initial polarized
factor. HD-R12 applies at its final rank because the exact
character field was retained at each step. The four forms
are the stated remaining construction tasks; the argument
provides no cycle for any of them.
